\chapter*{Ход выполнения практики}
\addcontentsline{toc}{chapter}{Ход выполнения практики}

\section*{Задания и отметки о выполнении}

\task{Написать обзор по предметной области исследования,
основываясь на научной, учебной и учебно-методической литературе, как на русском, так
и на английском языках. Необходимо использовать современную литературу на английском и русском языках по тематике научно-исследовательской работы, поиск которой можно осуществлять по библиографическим базам Scopus, WoS, elibrary, ResearchGate, ADS, ЭБС Лань и др.}

\task{Изучить особенности FDM-печати и типы возможных дефектов, возникающих в процессе.}

\task{Разработать удобный и интуитивно понятный пользовательский интерфейс приложения.}

\task{Обучить нейронную сеть распознаванию дефектов 3D-печати на основе изображений изделий.}

\task{Реализовать функционал приложения, включая загрузку данных, их обработку и визуализацию результатов.}

\task{Подготовить отчет по практике в соответствии с компетенциями направления 09.04.01 Информатика и вычислительная техника.}

\task{Создать презентацию в соответствии с данным отчетом и устный доклад для публичной защиты.}

\clearpage
\section*{Отзыв о работе обучающегося}

В ходе прохождения производственной практики, научно-исследовательской работы Ешенко Даниил Ярославович выполнил поставленные ему задачи, а именно изучил особенности FDM-печати и типы возможных дефектов, возникающих в её процессе, спроектировал интуитивно понятный пользовательский интерфейс мобильного приложения, а также обучил нейронную сеть распознаванию дефектов 3D-печати и реализовал функционал приложения, включая загрузку данных, их обработку и визуализацию результатов. По итогам работы было создано мобильное приложение на базе Android для обнаружения дефектов 3D-печати. Считаю, что его работа заслуживает положительной оценки.
\progressapprov